\documentclass[11pt]{article}
\usepackage[margin=1in]{geometry}
\usepackage{amsmath,amssymb,amsthm}
\usepackage{enumitem}
\usepackage{hyperref}

\newtheorem{theorem}{Theorem}
\newtheorem{lemma}{Lemma}
\newtheorem{corollary}{Corollary}

\newcommand{\otpi}{\widehat{\otimes}_{\pi}}
\newcommand{\weakstar}{w^*}
\newcommand{\CH}{\mathsf{CH}}
\newcommand{\Srange}{\mathcal{S}}
\newcommand{\Lop}{\mathcal{L}}

\title{A CH Counterexample to the WLD--Dunford--Pettis Projective Tensor Question}
\author{Run \texttt{fa\_banach\_001}, agent \texttt{agent\_lane\_10}}
\date{2026-06-18}

\begin{document}
\maketitle

\section*{Status}

This is a candidate counterexample under \(\CH\) to Question 8(b) of Aviles,
Martinez-Cervantes, Rodriguez, and Rueda Zoca, \emph{Topological properties in
tensor products of Banach spaces}, arXiv:2202.00371.

\section*{Source Question}

Question 8(b) asks:

\begin{quote}
Suppose that \(X\) and \(Y\) are WLD and that either \(X\) or \(Y\) has the
Dunford-Pettis property. Is \(X\widehat{\otimes}_{\pi}Y\) WLD?
\end{quote}

\section*{Claim}

\begin{theorem}
Assume \(\CH\).  Then there exist WLD Banach spaces \(X\) and \(Y\), with
\(X\) having the Dunford-Pettis property, such that
\[
X\otpi Y
\]
is not WLD.  In particular, under \(\CH\), Question 8(b) has a negative
answer.
\end{theorem}

\section*{Idea of the Proof}

The source paper already isolates exactly what can go wrong.  Its Theorem 4.1
says that \(X\otpi Y\) is WLD if and only if \(X\) and \(Y\) are WLD and every
operator in both cross directions
\[
X\to Y^*, \qquad Y\to X^*
\]
has norm-separable range.

Now take a WLD space \(Y\) whose dual ball fails property (M).  Under \(\CH\),
such a \(Y\) exists by the Plichko/Plebanek construction quoted in the source
paper.  The source paper's Lebesgue-Bochner corollary says that, for WLD
\(Y\), property (M) of \((B_{Y^*},\weakstar)\) is equivalent to
\(L_1(\mu,Y)\) being WLD for every finite measure \(\mu\).  Hence failure of
property (M) supplies a finite measure \(\mu\) for which \(L_1(\mu,Y)\) is not
WLD.  Since \(L_1(\mu,Y)\cong L_1(\mu)\otpi Y\), and \(L_1(\mu)\) is WCG
hence WLD and has the Dunford-Pettis property, this is the desired
counterexample.

\section*{Proof}

We use three facts recorded in the source paper.

\begin{enumerate}[label=(\arabic*)]
\item The source paper's Theorem 4.1 states that \(X\otpi Y\) is WLD if and
only if \(X\) and \(Y\) are WLD and
\[
\Lop(X,Y^*)=\Srange(X,Y^*), \qquad
\Lop(Y,X^*)=\Srange(Y,X^*),
\]
where \(\Srange\) denotes the operators with norm-separable range.

\item The source paper's Theorem 4.4 and Corollary 4.5 say that, for a Banach
space \(Y\), the following are equivalent:
\[
Y\text{ is WLD and }(B_{Y^*},\weakstar)\text{ has property (M)}
\]
and
\[
L_1(\mu,Y)\text{ is WLD for every finite measure }\mu.
\]

\item In its preliminaries on measures on Corson compact spaces, the source
paper recalls that under \(\CH\) there exist WLD Banach spaces \(Y\) for which
\((B_{Y^*},\weakstar)\) fails property (M), citing Plebanek's Corollary 4.4.
\end{enumerate}

Assume \(\CH\), and choose such a WLD Banach space \(Y\).  By (2), since
\((B_{Y^*},\weakstar)\) fails property (M), it is not true that
\(L_1(\mu,Y)\) is WLD for every finite measure \(\mu\).  Therefore there is a
finite measure \(\mu\) such that \(L_1(\mu,Y)\) is not WLD.

Set
\[
X:=L_1(\mu).
\]
Since \(\mu\) is finite, \(X\) is weakly compactly generated.  Indeed,
\[
\{f\in L_1(\mu): |f|\le 1\}
\]
is weakly compact by uniform integrability, and the linear span of this order
interval contains \(L_\infty(\mu)\), which is dense in \(L_1(\mu)\).  Hence
\(X\) is WCG and therefore WLD.  Also, \(L_1(\mu)\) has the Dunford-Pettis
property by the classical Dunford-Pettis theorem.

Finally, the standard Bochner-space identification gives an isometric
isomorphism
\[
L_1(\mu,Y)\cong L_1(\mu)\otpi Y=X\otpi Y.
\]
By the choice of \(\mu\), this space is not WLD.  Thus \(X\) and \(Y\) are WLD,
one factor \(X\) has the Dunford-Pettis property, but \(X\otpi Y\) is not WLD.

Equivalently, applying (1), the same example supplies an operator
\[
L_1(\mu)\to Y^*
\]
with nonseparable range.  This is precisely the obstruction forbidden by the
positive conclusion of Question 8(b).

\section*{Relation to the Existing Partial Result}

The earlier packet
\[
\texttt{solutions/partial/2202.00371\_c0\_projective\_tensor\_wld\_subcase/}
\]
proves in ZFC that \(c_0(\Gamma)\otpi Y\) is WLD for every WLD \(Y\).  The
present packet is complementary: it shows that under \(\CH\), the arbitrary
WLD+Dunford-Pettis factor case has a negative answer.  It does not invalidate
the \(c_0(\Gamma)\) positive subcase.

\section*{Novelty and Literature Check}

The run's cheap indexes contained the previous \(c_0(\Gamma)\) partial packet
for arXiv:2202.00371 but no packet for this \(L_1(\mu)\)/CH counterexample.

Local source searches and web searches were run for the exact source question
and close phrases including:
\[
\text{``Question 8(b)'' + ``Topological properties in tensor products of Banach spaces''},
\]
\[
\text{``WLD'' + ``Dunford-Pettis'' + ``projective tensor product'' + ``property (M)''},
\]
and
\[
\text{``Lebesgue-Bochner'' + ``WLD'' + ``property (M)''}.
\]
They found the source arXiv page and adjacent tensor-product/Dunford-Pettis
literature, but no later direct formulation of this CH counterexample.

This packet should therefore be treated as a short CH consequence of
the source paper's own characterizations together with the CH example quoted
from Plebanek, rather than as an already-explicit literature answer.

\section*{Limitations}

This is a counterexample under \(\CH\), not a ZFC counterexample.  The ZFC
status of Question 8(b) remains separate.  In models where all Corson compact
spaces have property (M), this particular construction is unavailable.

\section*{References}

\begin{enumerate}[label={[\arabic*]}]
\item A. Aviles, G. Martinez-Cervantes, J. Rodriguez, A. Rueda Zoca,
\emph{Topological properties in tensor products of Banach spaces},
arXiv:2202.00371.
\item G. Plebanek, \emph{Convex Corson compacta and Radon measures},
Fund. Math. \textbf{175} (2002), no. 2, 143--154.
\item J. Diestel, \emph{A survey of results related to the Dunford-Pettis
property}, Contemporary Mathematics \textbf{2} (1980), 15--60.
\end{enumerate}

\section*{Human Review Recommendation}

Review as a likely valid counterexample under \(\CH\).  The most
important audit points are:
\begin{enumerate}[label=(\roman*)]
\item the source Corollary 4.5 quantifier ``for every finite measure \(\mu\)'';
\item the CH existence of a WLD \(Y\) whose dual ball fails property (M), as
quoted in the source preliminaries;
\item the standard facts that finite-measure \(L_1(\mu)\) spaces are WCG/WLD
and have the Dunford-Pettis property.
\end{enumerate}

\end{document}
